% ==========================================================
% 논문 메타데이터 — 이 파일만 수정하면 모든 별표에 반영됩니다.
% ==========================================================

% ---- 학위 종류 ----
\newcommand{\DegreeKR}{석사}           % "석사" 또는 "박사"
\newcommand{\FieldKR}{공학}            % "공학" / "이학" / "문학" 등

% ---- 논문 제목 ----
\newcommand{\ThesisTitleKR}{가천대학교 학위논문 \LaTeX\ 템플릿 샘플}
\newcommand{\ThesisTitleEN}{A Sample \LaTeX\ Template for Gachon University Thesis}

% ---- 저자 ----
% 표지용 성명은 한글 한 글자씩 공백을 두는 관례가 있습니다.
\newcommand{\AuthorKR}{홍\hspace{1.2em}길\hspace{1.2em}동}
\newcommand{\AuthorKRPlain}{홍길동}    % 인준서/초록 본문용 (공백 없는 성명)
\newcommand{\AuthorEN}{Gildong Hong}

% ---- 소속 ----
\newcommand{\DepartmentKR}{설비·소방공학과}
\newcommand{\DepartmentEN}{Department of Fire and Disaster Prevention Engineering}
\newcommand{\MajorKR}{소방방재공학 전공}
\newcommand{\GraduateSchoolKR}{가천대학교 대학원}
\newcommand{\GraduateSchoolEN}{Graduate School of Gachon University}

% ---- 지도교수 ----
\newcommand{\Advisor}{○\hspace{0.8em}○\hspace{0.8em}○}

% ---- 제출 시점 ----
\newcommand{\SubmissionYear}{2026}
\newcommand{\SubmissionMonth}{2}       % 2 또는 8

% ---- 심사위원 (인준서에서 사용) ----
% 석사: 3인, 박사: 5인
% 박사의 경우 frontmatter/04_approval.tex 에서 추가 2줄을 주석 해제하세요.
